\chapter{System Architecture / Arquitetura do Sistema}

\section*{English}
The network is hierarchical. Local clusters with baselines of tens of kilometres serve regional illuminators and low-altitude scenarios. Regional clusters extend coverage, while long baselines may support HF, high-altitude and space-object observations when common visibility and a compatible illuminator exist.

Each station contains a radome platform, multiband RF chains, local processing, timing and calibration, optical transport, power conditioning, thermal management and a supervisory interface. The central fusion layer associates detections by waveform, time, frequency, Doppler, polarization and geometric consistency. It must distinguish direct emitter observations, bistatic reflections, interference and tracking hypotheses.

\begin{figure}[H]
\centering
\includegraphics[width=0.95\textwidth]{fig01_arquitetura_rede.png}
\caption{Distributed sensing architecture / Arquitetura distribuída de sensoriamento.}
\label{fig:architecture}
\end{figure}

\section*{Português}
A rede possui arquitetura hierárquica. Clusters locais, com linhas de base de dezenas de quilômetros, atendem iluminadores regionais e cenários de baixa altitude. Clusters regionais ampliam a cobertura, enquanto linhas de base longas podem apoiar observações em HF, alvos de alta altitude e objetos espaciais quando há visada comum e iluminador compatível.

Cada estação contém plataforma radome, cadeias RF multifaixa, processamento local, sincronismo e calibração, transporte óptico, condicionamento de energia, gestão térmica e interface de supervisão. A fusão central associa detecções por forma de onda, tempo, frequência, Doppler, polarização e consistência geométrica. Deve distinguir observações diretas de emissores, reflexões biestáticas, interferência e hipóteses de rastreamento.
